% !TeX root = ../navier-stokes-manuscript.tex
\subsection{The critical-height problem inside a stretching vortex}\label{sub:ThePerfectlyStretchingVortex}

Follow the maximal history into a region where the fluid is rotating
coherently and the surrounding velocity field is pulling along the direction
of that rotation.  The two local objects are
\[
  \omega:=\nabla\times u,
  \qquad
  S:=\frac12\bigl(\nabla u+(\nabla u)^{\mathsf T}\bigr).
\]
Taking the curl removes the pressure gradient but not the motion through which
pressure helped create the velocity field.  The rotating part of that same
field obeys
\begin{equation}\label{eq:BodyVorticityEquation}
  \partial_t\omega+(u\cdot\nabla)\omega
  =S\omega+\nu\Delta\omega.
\end{equation}
The right-hand side contains the entire contest.  Viscosity diffuses variation
in the vorticity.  The term \(S\omega\) is the action of the local strain on
the direction and magnitude of the rotation.  When \(\omega\) points along an
extensional direction of \(S\),
\[
  \omega\cdot S\omega>0,
\]
and the surrounding motion feeds the rotation rather than flattening it.

It cannot receive that amplification while keeping the same shape.
Incompressibility says
\[
  \operatorname{tr}S=\nabla\cdot u=0.
\]
so extension along the vortex axis must be paid for by contraction across the
axis.  The same rotating structure becomes longer and thinner.  A velocity
variation \(\Delta U\) that once occupied a transverse width \(\ell\) is now
carried through less space, and its gradient has the scale
\[
  |\nabla u|\sim\frac{\Delta U}{\ell}.
\]
The numerator need not become infinite for the ratio to do so.  A finite
velocity difference can become an unbounded derivative if the distance across
which the fluid makes that difference is allowed to collapse.

Nothing in the kinetic-energy balance forbids this geometry.  For the same
solution, that balance is
\begin{equation}\label{eq:BodyEnergyIdentity}
  \frac12\norm{u(t)}_2^2
  +\nu\int_0^t\norm{\nabla u(s)}_2^2\,ds
  =\frac12\norm{u_0}_2^2.
\end{equation}
The first term measures velocity across the volume of the fluid.  It does not
remember the width across which any one part of that velocity changes.  The
second term pays for the square of the gradient only after time has been
integrated.  A steep event may occupy a short interval, a steeper event a
shorter one, and the account may remain finite while the pointwise spatial
profile becomes harder and harder to resolve.

The narrowing also strengthens viscosity, so stretching by itself has not
identified a route to singularity.  Pair
\eqref{eq:BodyVorticityEquation} with \(\omega\).  Before integrating through
space, the rotation satisfies
\[
  \frac12(\partial_t+u\cdot\nabla)|\omega|^2
  =\omega\cdot S\omega
  +\frac\nu2\Delta|\omega|^2
  -\nu|\nabla\omega|^2.
\]
After integration over \(\R^3\), transport carries enstrophy through the fluid
without changing its total amount, and the Laplacian flux vanishes at spatial
infinity.  What remains is the exact competition
\begin{equation}\label{eq:BodyEnstrophyBalance}
  \frac12\frac d{dt}\norm{\omega(t)}_2^2
  +\nu\norm{\nabla\omega(t)}_2^2
  =\int_{\R^3}\omega\cdot S\omega\,dx.
\end{equation}
The term on the right can produce enstrophy; the term on the left removes it.
Both act on the same rotating structure and both become stronger as its width
falls.  Their magnitudes at one chosen width cannot tell us which action gains
when the whole structure is made smaller.  For that we have to shrink the law
itself and watch how every term changes together.

So shrink the entire equation without changing the law that governs it.  For
\(\lambda>0\), define
\begin{equation}\label{eq:BodyNavierStokesScaling}
  u_\lambda(x,t):=\lambda u(\lambda x,\lambda^2t),
  \qquad
  p_\lambda(x,t):=\lambda^2p(\lambda x,\lambda^2t).
\end{equation}
The time response, transport, pressure response, and viscous response all
receive the same factor \(\lambda^3\):
\[
\begin{aligned}
  \partial_tu_\lambda
  &=\lambda^3(\partial_tu)(\lambda x,\lambda^2t),
  &
  (u_\lambda\cdot\nabla)u_\lambda
  &=\lambda^3((u\cdot\nabla)u)(\lambda x,\lambda^2t),
  \\
  \nabla p_\lambda
  &=\lambda^3(\nabla p)(\lambda x,\lambda^2t),
  &
  \Delta u_\lambda
  &=\lambda^3(\Delta u)(\lambda x,\lambda^2t).
\end{aligned}
\]
The smaller copy has not handed viscosity an extra power with which to defeat
transport.  The terms remain locked in the proportion already fixed by the
equation.  What has changed is the room and the time: space contracts by
\(\lambda^{-1}\), and the same lawful motion is run on a clock shorter by
\(\lambda^{-2}\).

We now need the spatial reading that does not fade when the structure shrinks
and does not inflate simply because we have magnified it.  Let
\(\Lambda=(-\Delta)^{1/2}\).  At Sobolev height \(s\), a change of variables
gives
\begin{equation}\label{eq:BodyHomogeneousSobolevScaling}
  \norm{u_\lambda(t)}_{\dot H^s}^2
  =\lambda^{2s-1}\norm{u(\lambda^2t)}_{\dot H^s}^2.
\end{equation}
There is one height at which the scale factor disappears:
\[
  2s-1=0,
  \qquad
  s=\frac12.
\]
This selects
\begin{equation}\label{eq:BodyCriticalHeightDefinition}
  H(t):=\frac12\norm{u(t)}_{\dot H^{1/2}}^2
  =\frac12\norm{\Lambda^{1/2}u(t)}_2^2
\end{equation}
as the critical height.  A lawful concentration cannot lower this reading by
making itself smaller.  Kinetic energy sits below it and loses one power of
\(\lambda\); the first derivative sits above it and gains one:
\[
  \norm{u_\lambda(t)}_2^2
  =\lambda^{-1}\norm{u(\lambda^2t)}_2^2,
  \qquad
  H_\lambda(t)=H(\lambda^2t),
  \qquad
  \norm{\nabla u_\lambda(t)}_2^2
  =\lambda\norm{\nabla u(\lambda^2t)}_2^2.
\]
The same contraction can now make energy quieter and a derivative louder.
Between them, the critical height keeps the size of the concentration fixed.
It is the scale at which the narrowing can no longer hide behind the energy
balance.

One rescaled profile lets us watch all three readings at once.  For a nonzero
divergence-free \(\phi\in\Ssigma\), set
\[
  u_\ell(x):=\ell^{-1}\phi(x/\ell).
\]
Then
\[
\begin{aligned}
  \norm{u_\ell}_\infty&\sim\ell^{-1},
  &\norm{\nabla u_\ell}_\infty&\sim\ell^{-2},
  \\
  \norm{u_\ell}_2^2&\sim\ell,
  &\norm{u_\ell}_{\dot H^{1/2}}^2
  &=\norm{\phi}_{\dot H^{1/2}}^2.
\end{aligned}
\]
As \(\ell\downarrow0\), the amplitude and gradient rise, the energy falls, and
the critical height does not move.  These profiles are not successive states
of a Navier--Stokes solution; nothing here says that the evolution can produce
them.  They isolate the shape of the loophole: energy allows it, scaling
preserves it, and the nonlinear equation still has to decide whether the
actual fluid can travel through it.

Scaling has selected the height at which a shrinking structure remains fully
visible.  We can now return to the equation and ask how the actual solution
moves through that height.  For
\[
  E_s(t):=\frac12\norm{\Lambda^su(t)}_2^2,
\]
keep the transport and pressure generated by the same velocity together in
the signed response
\[
  \mathcal P_s(t)
  :=-
  \operatorname{Re}
  \pair{\Lambda^su}
  {\Lambda^s\bigl((u\cdot\nabla)u+\nabla p\bigr)}_{L^2}.
\]
At this global scalar level the pressure contribution vanishes after
incompressibility is used.  It remains in the definition because the balance
was produced by the pressure-complete momentum equation, and the differentiated
rows later in the paper must remember that origin.  Pairing at height \(s\)
gives
\begin{equation}\label{eq:BodySpatialEnergyBalanceIntro}
  E_s'=\mathcal P_s-2\nu E_{s+1}.
\end{equation}
At the critical height,
\begin{equation}\label{eq:BodyCriticalHeightBalanceIntro}
  H'=\mathcal P_{1/2}-2\nu E_{3/2}.
\end{equation}
Thus a positive nonlinear production can raise \(H\) only against the viscous
payment \(2\nu E_{3/2}\) made at the same instant.  The geometric contest in
the vortex has become one exact scalar row.  Neither stretching nor diffusion
can win this row by changing scale; the unresolved issue is whether the signed
production can keep outrunning its adjacent viscous cost along the actual
history.

Now suppose, only to follow the possible failure, that the maximal history has
a finite endpoint \(T_*\).  The endpoint criterion of Escauriaza, Seregin, and
\v{S}ver{\'a}k says that a uniform \(L^3\) ceiling would continue the solution
through that endpoint \parencite{EscauriazaSereginSverak2003}.  The critical
Sobolev embedding reads
\begin{equation}\label{eq:BodyCriticalHeightControlsLThree}
  \norm{u(t)}_3
  \leq C\norm{\Lambda^{1/2}u(t)}_2
  =C\sqrt{2H(t)}.
\end{equation}
Consequently the same hypothetical endpoint forces
\begin{equation}\label{eq:BodyFiniteEndpointForcesCriticalEscape}
  \sup_{0\leq t<T_*}H(t)=\infty.
\end{equation}
This is no longer a sequence of profiles chosen outside the evolution.  The
actual velocity history must climb through arbitrarily large critical levels.
If
\[
  M_n:=2^n\max\{1,H(0)\},
  \qquad
  t_n:=\inf\{t<T_*:H(t)=M_n\},
\]
then continuity gives \(t_n\uparrow T_*\).  The first-passage intervals are
successive pieces of one history, and
\begin{equation}\label{eq:BodyActualZenoIntervals}
  \sum_{n=0}^{\infty}(t_{n+1}-t_n)=T_*-t_0<\infty.
\end{equation}

The heat equation now gives this finite-time ascent its spatial clock.  If
viscosity is acting on a structure of width \(r\), its linear comparison time
is determined without analogy.  For \(v_t=\nu\Delta v\),
\[
  \widehat v(\xi,t+\delta t)
  =e^{-\nu|\xi|^2\delta t}\widehat v(\xi,t).
\]
A structure of width \(r\) lives near \(|\xi|\simeq r^{-1}\), so diffusion
changes its amplitude by order one on the time
\begin{equation}\label{eq:BodyHeatComparisonTime}
  \tau_\nu(r)\asymp\frac{r^2}{\nu}.
\end{equation}
This does not say that the nonlinear fluid forms that width, or that one width
can be read uniquely from \(H\).  It says what time viscosity assigns to a
scale once the evolution has formed it.

Take geometrically shrinking widths
\[
  r_n:=2^{-n}r_0,
  \qquad
  \tau_n:=\frac{r_n^2}{\nu}.
\]
Their heat clocks satisfy
\begin{equation}\label{eq:BodyZenoHeatClock}
  \tau_n=\frac{r_0^2}{\nu}4^{-n},
  \qquad
  \sum_{n=0}^{\infty}\tau_n
  =\frac{4r_0^2}{3\nu}<\infty.
\end{equation}
Thus the viscous clocks attached to infinitely many geometrically smaller
scales fit inside a finite amount of elapsed time.  The arithmetic agrees with
the actual first-passage compression in
\eqref{eq:BodyActualZenoIntervals}; it does not identify those passages with
the model widths \(r_n\).

For a singularity to use this opening, the fluid would have to form
\(r_0\), then \(r_1\), then \(r_2\), without losing the stretching alignment
while viscosity acts on every scale already formed.  Each event would begin
from the velocity produced by the previous event.  The widths would shrink,
the clocks would shorten, and the critical production in
\eqref{eq:BodyCriticalHeightBalanceIntro} would have to meet its viscous cost
again and again.

The spatial picture has now become a temporal obligation.  The same velocity
must cross unbounded critical levels during intervals whose total length is
finite.  Those changes have to be carried by its time response; if that
response also escapes, its changes have to be carried by the next one.  This
is where the temporal derivative Tower begins.
